\begin{abstract}

Ionic thermoelectric materials provide a route for converting low-grade heat
into electricity through thermodiffusion or thermogalvanic processes. Their
performance is governed by ionic thermodynamics and transport, particularly
the reaction entropy and solvation environments of redox couples. This thesis
examines the molecular design of thermogalvanic electrolytes, with emphasis on
solvent effects and polymer-confined hydrogel systems. For liquid
ferricyanide--ferrocyanide electrolytes, the limited availability of
thermopower data and the complexity of ion--solvent interactions make direct
machine-learning prediction difficult. Density functional theory calculations
were therefore combined with a small-data machine-learning approach for
solvent screening. The resulting analysis identified solvent polarity as an
important factor governing thermopower. This relationship was further examined
in poly(vinyl alcohol) hydrogels. A
volcano-shaped dependence of thermopower on solvent dipole moment was observed,
and fine control of hydrogel polarity produced a thermopower of
3.28 mV K$^{-1}$. The enhancement was associated with asymmetric changes in
the solvation environments of the redox ions. These findings provide a basis
for relating solvent properties and redox-ion solvation to the performance of
thermogalvanic electrolytes.

\end{abstract}
